\diarytitle{0011}{一次同次線形偏微分方程式の一般解（$x u_x - y u_y = 0$）}

$P(x,y)\,\dfrac{\partial u}{\partial x} + Q(x,y)\,\dfrac{\partial u}{\partial y} = 0$ の一般解は、特性方程式 $\dfrac{dx}{P} = \dfrac{dy}{Q}$ の解 $\varphi(x,y) = C$ を用いて $u = F(\varphi(x,y))$（$F$ は任意関数）と表せる。

特性方程式は
\[
\frac{dx}{x} = \frac{dy}{-y}
\]
両辺を積分して $\ln|x| = -\ln|y| + C_{1}$、すなわち
\[
xy = C
\]
が特性曲線である。したがって一般解は、任意関数 $F$ を用いて
\[
\boxed{\; u(x,y) = F(xy) \;}
\]
（検算: $u_x = y\,F'(xy)$, $u_y = x\,F'(xy)$ より $x u_x - y u_y = xy\,F' - xy\,F' = 0$。）
